\documentclass[12pt]{article}
\usepackage[T1]{fontenc}
\usepackage[utf8]{inputenc}
\usepackage{lmodern}
\usepackage{textcomp}
\usepackage{enumitem}
\usepackage{geometry}
\geometry{margin=1in}
\begin{document}
\begin{center}
Answer Key - Constitutional Law - Graduate Level Assessment 25 \
\small (Answers are ordered exactly as in the question paper.)
\end{center}

\section*{Multiple Choice}
\begin{enumerate}[label=\arabic*.]
\item \textbf{Answer:} D
\\ \textit{Options:} A) Collateral matters; B) Public records; C) When it's not offered to prove the contents of the document; D) All of the above.
\\ \textit{Note:} The best evidence rule generally does not apply in situations involving public records, summaries of voluminous documents, or when the original document is not required for authenticity, which encompasses all the options presented in the question.
\item \textbf{Answer:} A
\\ \textit{Options:} A) It is an agreement between two or more to commit a crime.; B) It is solicitation of another to commit a felony.; C) No act is needed other than the solicitation.; D) Withdrawal is generally not a defense.
\\ \textit{Note:} The correct answer is that solicitation does not involve an agreement to commit a crime, as solicitation specifically refers to the act of encouraging or requesting another person to engage in criminal conduct, rather than forming a conspiratorial agreement.
\item \textbf{Answer:} C
\\ \textit{Options:} A) Preliminary Hearing; B) Arraignment; C) investigative surveillance; D) Post-charge lineups
\\ \textit{Note:} The correct answer is investigative surveillance because the Sixth Amendment right to counsel does not attach until formal charges have been filed, and at the surveillance stage, the individual is not yet subject to prosecution.
\item \textbf{Answer:} A
\\ \textit{Options:} A) not hearsay.; B) hearsay, but admissible as an admission.; C) hearsay, but admissible as a dying declaration.; D) hearsay not within any recognized exception.
\\ \textit{Note:} The testimony is not hearsay because it is offered for the truth of the matter asserted-that someone else committed the murder-rather than to prove the truth of the statement made by the witness on death row.
\item \textbf{Answer:} C
\\ \textit{Options:} A) A breaking and entry; B) Of a dwelling, of another; C) During the daytime; D) With the intent to commit a felony therein.
\\ \textit{Note:} "During the daytime" is correct because common law burglary specifically requires the act to occur at night, distinguishing it from other forms of unlawful entry.
\end{enumerate}

\section*{True / False}
\begin{enumerate}[label=\arabic*.]
\setcounter{enumi}{5}
\item \textbf{Answer:} True
\\ \textit{Options:} A) True; B) False
\\ \textit{Note:} Cold pursuit is not recognized as a valid exception to the warrant requirement because the legal principle of exigent circumstances necessitates an immediate need for law enforcement to act, which is not present in situations where there is time to obtain a warrant.
\item \textbf{Answer:} True
\\ \textit{Options:} A) True; B) False
\\ \textit{Note:} The Doctrine of Worthier Title holds that a remainder to the grantor's heirs is invalid because it contradicts the principle that a grantor cannot create a future interest in their own heirs, thereby resulting in a reversion of the property back to the grantor.
\item \textbf{Answer:} False
\\ \textit{Options:} A) True; B) False
\\ \textit{Note:} The statement is false because the U.S. Constitution requires that Supreme Court justices be appointed by the President with the advice and consent of the Senate, as outlined in Article II, Section 2.
\item \textbf{Answer:} False
\\ \textit{Options:} A) True; B) False
\\ \textit{Note:} The enabling clause of the Fourteenth Amendment primarily empowers Congress to enforce the rights guaranteed by the amendment against state actions, rather than allowing states to regulate private individuals' actions based on nationality.
\item \textbf{Answer:} False
\\ \textit{Options:} A) True; B) False
\\ \textit{Note:} Slander per se does not require publication to a third party because it is actionable based on the nature of the statement itself, which is considered inherently harmful, regardless of its dissemination.
\end{enumerate}

\section*{Long Form Response}
\begin{enumerate}[label=\arabic*.]
\setcounter{enumi}{10}
\item \textbf{Answer:} The necessary and proper power, articulated in Article I, Section 8 of the U.S. Constitution, allows Congress to enact laws essential for executing its enumerated powers. However, this power cannot stand alone to support federal authority because it is inherently dependent on the existence of other constitutional provisions that delegate specific powers to the federal government. The necessary and proper clause serves as a conduit for implementing those enumerated powers, rather than a source of independent authority. Thus, without clearly defined powers, the necessary and proper clause risks overreach and undermines the balance of federalism by allowing Congress to legislate beyond its intended scope. This relationship underscores the importance of a constitutional framework that maintains limits on federal authority while enabling effective governance.
\\ \textit{Note:} correct because the necessary and proper power functions to facilitate the execution of Congress's enumerated powers, relying on the existence of those specific powers to justify federal authority, rather than serving as an independent source of power.
\item \textbf{Answer:} Denying certain fundamental rights to specific individuals while granting them to others fundamentally undermines the principle of equal protection under the law, as enshrined in the Fourteenth Amendment of the U.S. Constitution. This principle mandates that individuals in similar circumstances be treated equally, thereby fostering a sense of justice and societal cohesion. When the law discriminates, it not only inflicts harm on the marginalized groups but also erodes public trust in legal institutions, leading to societal fragmentation. Such disparities in rights can perpetuate systemic inequalities and hinder the pursuit of a truly democratic society, as they create hierarchies of worth among citizens. Ultimately, the failure to uphold equal protection not only contravenes constitutional mandates but also compromises the foundational ideals of liberty and justice for all.
\\ \textit{Note:} correct because it articulates that denying fundamental rights to certain individuals violates the equal protection clause by treating similarly situated individuals unequally, which harms marginalized groups and diminishes public trust in the legal system.
\end{enumerate}

\vfill
\noindent\textit{End of Answer Key}
\end{document}